% ==============================================================================
% Fichier     : chapitres/Apropos_Auteur.tex
% Titre       : À Propos de l'Auteur
% Auteur      : MINKA MI NGUIDJOÏ Thierry Emmanuel
% ==============================================================================

\chapter*{À Propos de l'Auteur}
\addcontentsline{toc}{chapter}{À Propos de l'Auteur}
\markboth{À Propos de l'Auteur}{À Propos de l'Auteur}

\vspace{1ex}

% Bandeau décoratif supérieur
\begin{tcolorbox}[
    enhanced,
    colback=CouleurPrimaire,
    colframe=CouleurPrimaire,
    arc=0pt,
    boxrule=0pt,
    left=0pt, right=0pt, top=4pt, bottom=4pt
]
\centering
{\color{CouleurAccent}\rule{0.3\linewidth}{0.8pt}}\quad
{\color{white}\large\bfseries MINKA MI NGUIDJOÏ Thierry Emmanuel}\quad
{\color{CouleurAccent}\rule{0.3\linewidth}{0.8pt}}\\[2pt]
{\color{white!80!gray}\small\itshape
    Chercheur~$\cdot$~Expert en Sécurité de l'Information~$\cdot$~Enseignant}
\end{tcolorbox}

\vspace{1.5ex}

% ==============================================================================
% BLOC BIOGRAPHIE PRINCIPALE
% ==============================================================================

\begin{tcolorbox}[
    enhanced,
    breakable,
    colback=white,
    colframe=CouleurSecondaire,
    arc=4pt,
    boxrule=1pt,
    left=10pt, right=10pt, top=10pt, bottom=10pt
]

\textbf{MINKA MI NGUIDJOÏ Thierry Emmanuel} est chercheur au \termecle{Laboratoire d'Informatique et de Mathématiques pour les Sciences de l'Ingénieur} (\sigle{LIMSI}) de l'Université de Yaoundé~I, au Cameroun. Ses travaux de doctorat, conduits sous la double direction du Professeur \textbf{Flavien Serge Mani Onana} et du Professeur \textbf{Thomas Djotio Ndié}, portent sur les fondements formels de la \termecle{cryptographie institutionnelle} — une discipline qu'il a lui-même fondée et dont il est l'initiateur reconnu au sein de la communauté scientifique internationale.

\medskip

Sa thèse, intitulée \textit{«~Aux fondements de la cryptographie institutionnelle~: formalisation, bornes d'impossibilité, conditions d'optimalité et vérification formelle~»}, articule des contributions théoriques majeures autour du \textbf{Cadre CLO} (\textit{Confidentiality-Legibility-Opacity}), du \textbf{Trilemme CRO} (\textit{Confidentiality-Recoverability-Opacity}), de la \textbf{Famille SOTP} (\textit{State-based One-Time Protocols}) et des extensions post-quantiques qui en découlent. Ces travaux, formalisés en partie avec l'assistant de preuve \textbf{Lean 4}, ont donné lieu à neuf prépublications sur l'\sigle{IACR ePrint Archive} et à plusieurs articles publiés dans \textit{IEEE Access}.

\end{tcolorbox}

\vspace{1ex}

% ==============================================================================
% DEUX COLONNES : PARCOURS ET CERTIFICATIONS
% ==============================================================================

\begin{minipage}[t]{0.52\textwidth}
\begin{tcolorbox}[
    enhanced,
    colback=CouleurDef,
    colframe=CouleurBordDef,
    fonttitle=\bfseries\color{white},
    title={\textbf{[\S]}\quad Parcours Académique et Institutionnel},
    attach boxed title to top left={yshift=-2mm, xshift=4mm},
    boxed title style={colback=CouleurBordDef, rounded corners},
    arc=3pt, boxrule=0.8pt,
    left=6pt, right=6pt, top=8pt, bottom=6pt
]
\begin{itemize}
    \item \textbf{Doctorant} en Informatique\\
          {\footnotesize Université de Yaoundé~I — \sigle{LIMSI}\\
          \textit{Soutenance prévue : juin 2026}}
    \item \textbf{Chercheur associé} en cryptographie\\
          {\footnotesize Formalisation sous Lean 4, preuves post-quantiques}
    \item \textbf{Enseignant} --- Investigation Numérique\\
          {\footnotesize Cours dispensé en cycle d'ingénierie informatique}
    \item \textbf{Responsable SSI} --- Projet GIP Bénin\\
          {\footnotesize Gouvernance et sécurité des SI publics}
    \item \textbf{Cadre supérieur} --- Contrôle Supérieur de l'État\\
          {\footnotesize Institution supérieure d'audit — République du Cameroun}
\end{itemize}
\end{tcolorbox}
\end{minipage}
\hfill
\begin{minipage}[t]{0.44\textwidth}
\begin{tcolorbox}[
    enhanced,
    colback=CouleurLizbiz,
    colframe=CouleurBordLizbiz,
    fonttitle=\bfseries\color{white},
    title={\ding{166}\quad Certifications Professionnelles},
    attach boxed title to top left={yshift=-2mm, xshift=4mm},
    boxed title style={colback=CouleurBordLizbiz, rounded corners},
    arc=3pt, boxrule=0.8pt,
    left=6pt, right=6pt, top=8pt, bottom=6pt
]
\begin{itemize}
    \item \textbf{\sigle{CISA}}\\
          {\footnotesize \textit{Certified Information Systems Auditor}}
    \item \textbf{\sigle{CISM}}\\
          {\footnotesize \textit{Certified Information Security Manager}}
    \item \textbf{\sigle{CGEIT}}\\
          {\footnotesize \textit{Certified in the Governance of Enterprise IT}}
    \item \textbf{\sigle{CRISC}}\\
          {\footnotesize \textit{Certified in Risk \& Information Systems Control}}
    \item \textbf{\sigle{CDPSE}}\\
          {\footnotesize \textit{Certified Data Privacy Solutions Engineer}}
\end{itemize}
{\footnotesize\color{CouleurBordInfo}\itshape
Toutes certifications délivrées par l'\textbf{ISACA} — l'organisation mondiale de référence en gouvernance, audit et sécurité des SI.}
\end{tcolorbox}
\end{minipage}

\vspace{1.5ex}

% ==============================================================================
% CONTRIBUTIONS SCIENTIFIQUES
% ==============================================================================

\begin{tcolorbox}[
    enhanced,
    colback=white,
    colframe=CouleurAccent,
    fonttitle=\bfseries\color{white},
    title={\ding{167}\quad Contributions Scientifiques Sélectionnées},
    attach boxed title to top left={yshift=-2mm, xshift=4mm},
    boxed title style={colback=CouleurAccent, rounded corners},
    arc=3pt, boxrule=0.8pt,
    left=8pt, right=8pt, top=10pt, bottom=8pt
]

\begin{itemize}
    \item \textbf{P1 — Cadre CLO} : \textit{A Formal Framework for Institutional Cryptography: The CLO Model}\\
          {\footnotesize IEEE Access, pp.~20427--20445, DOI:~\texttt{10.1109/ACCESS.2026.3660105}}

    \item \textbf{P2 — Trilemme CRO} : \textit{The CRO Trilemma: Impossibility Bounds in Institutional Cryptographic Protocols}\\
          {\footnotesize IEEE Access, pp.~33263--33274, DOI:~\texttt{10.1109/ACCESS.2026.3669065}}

    \item \textbf{S1 — Extension PQ-CRO} : Extension post-quantique du Trilemme CRO\\
          {\footnotesize \textit{Soumis à IEEE Access} --- MS-ID:~\texttt{Access-2026-09464}}

    \item \textbf{S2 — Famille SOTP} : \textit{The SOTP Protocol Family: State-based One-Time Protocols}\\
          {\footnotesize \textit{Soumis à IEEE Access} --- MS-ID:~\texttt{Access-2026-09012}}

    \item \textbf{Neuf prépublications} disponibles sur l'\textbf{IACR ePrint Archive}\\
          {\footnotesize Couvrant les primitives CEE, AOW, ZK-NR, Q2CSI, AIIP et les formalisations Lean~4}
\end{itemize}

\end{tcolorbox}

\vspace{1ex}

% ==============================================================================
% GENÈSE DE CE COURS
% ==============================================================================

\begin{tcolorbox}[
    enhanced,
    breakable,
    colback=CouleurAlerte,
    colframe=CouleurBordAlerte,
    fonttitle=\bfseries\color{white},
    title={\ding{96}\quad Genèse de ce Cours : D'une Expérience de Terrain à une Discipline},
    attach boxed title to top left={yshift=-2mm, xshift=4mm},
    boxed title style={colback=CouleurBordAlerte, rounded corners},
    arc=3pt, boxrule=0.8pt,
    left=8pt, right=8pt, top=10pt, bottom=8pt
]

Ce cours d'Audit des Systèmes d'Information trouve sa source dans une expérience concrète vécue par l'auteur entre 2018 et 2019, lors d'une mission d'audit du système de gestion de la paie \textbf{SIGIPES} au sein de l'administration camerounaise. Confronté à la tension fondamentale entre la \textit{confidentialité} des données personnelles des agents de l'État et l'\textit{admissibilité légale} des preuves numériques dans un contexte de contrôle institutionnel, l'auteur a mesuré les limites des paradigmes cryptographiques classiques — qui conçoivent la confidentialité comme une propriété absolue, indifférente aux exigences du droit et des institutions.

\medskip

C'est cette expérience d'audit qui a semé les graines de la \textbf{cryptographie institutionnelle} — une discipline désormais formalisée, publiée et portée devant la communauté scientifique internationale. Ce cours en est le miroir pédagogique : transmettre aux praticiens de demain la culture du contrôle, de la preuve et de la gouvernance numérique qui permet à l'audit de jouer pleinement son rôle de garant de la confiance organisationnelle.

\end{tcolorbox}

\vspace{1.5ex}

% ==============================================================================
% CONTACT
% ==============================================================================

\begin{center}
{\color{CouleurAccent}\rule{0.5\linewidth}{0.6pt}}\\[0.8ex]
{\small\color{CouleurPrimaire}
    \ding{229}~\textit{Université de Yaoundé~I, LIMSI — Yaoundé, Cameroun}\\[0.3ex]
    \ding{41}~\textit{Contacts disponibles auprès de l'institution}\\[0.3ex]
    \ding{232}~\textit{Publications : \texttt{https://eprint.iacr.org} (recherche : MINKA)}
}\\[0.8ex]
{\color{CouleurAccent}\rule{0.5\linewidth}{0.6pt}}
\end{center}
